% HAAP SDK — Binary Distribution & CrewAI Integration Guide
% Compile with: pdflatex DISTRIBUTION.tex  (twice for TOC)
% Requires: listings, booktabs, geometry, hyperref, xcolor, fancyhdr, mdframed

\documentclass[11pt, a4paper]{article}

% ── Page geometry ────────────────────────────────────────────────────────────
\usepackage[
  top=2.5cm, bottom=2.5cm,
  left=2.8cm, right=2.8cm,
  headheight=14pt
]{geometry}

% ── Fonts & encoding ─────────────────────────────────────────────────────────
\usepackage[T1]{fontenc}
\usepackage[utf8]{inputenc}
\usepackage{lmodern}
\usepackage{microtype}

% ── Colour ───────────────────────────────────────────────────────────────────
\usepackage[dvipsnames]{xcolor}
\definecolor{hawcxblue}{HTML}{1A56DB}
\definecolor{hawcxgray}{HTML}{F3F4F6}
\definecolor{codered}{HTML}{C0392B}
\definecolor{codegreen}{HTML}{1E8449}
\definecolor{codegray}{HTML}{7F8C8D}
\definecolor{codebg}{HTML}{F8F9FA}
\definecolor{frameborder}{HTML}{D1D5DB}

% ── Tables ───────────────────────────────────────────────────────────────────
\usepackage{booktabs}
\usepackage{array}
\usepackage{tabularx}

% ── Code listings ────────────────────────────────────────────────────────────
\usepackage{listings}
\lstset{
  backgroundcolor=\color{codebg},
  basicstyle=\ttfamily\small,
  breaklines=true,
  breakatwhitespace=true,
  captionpos=b,
  commentstyle=\color{codegray}\itshape,
  frame=single,
  framesep=4pt,
  framerule=0.4pt,
  rulecolor=\color{frameborder},
  keywordstyle=\color{hawcxblue}\bfseries,
  numberstyle=\tiny\color{codegray},
  numbers=left,
  numbersep=8pt,
  showstringspaces=false,
  stringstyle=\color{codered},
  tabsize=4,
  xleftmargin=1.5em,
  xrightmargin=0.5em,
}
\lstdefinelanguage{TOML}{
  morestring=[b]",
  morestring=[b]',
  morecomment=[l]{\#},
  keywords={true, false},
  sensitive=false,
}
\lstdefinelanguage{JSON}{
  morestring=[b]",
  literate=
    *{:}{{{\color{codegray}{:}}}}{1}
     {,}{{{\color{codegray}{,}}}}{1}
     {\{}{{{\color{codegray}{\{}}}}{1}
     {\}}{{{\color{codegray}{\}}}}}{1}
     {[}{{{\color{codegray}{[}}}}{1}
     {]}{{{\color{codegray}{]}}}}{1},
}

% ── Frames / boxes ───────────────────────────────────────────────────────────
\usepackage{mdframed}
\mdfdefinestyle{notebox}{
  backgroundcolor=hawcxgray,
  linecolor=hawcxblue,
  linewidth=2pt,
  leftline=true,
  rightline=false,
  topline=false,
  bottomline=false,
  innerleftmargin=10pt,
  innerrightmargin=10pt,
  innertopmargin=6pt,
  innerbottommargin=6pt,
}

% ── Hyperlinks & PDF metadata ────────────────────────────────────────────────
\usepackage[
  colorlinks=true,
  linkcolor=hawcxblue,
  urlcolor=hawcxblue,
  citecolor=hawcxblue,
  pdftitle={HAAP SDK — Binary Distribution \& CrewAI Integration Guide},
  pdfauthor={Hawcx Inc.},
  pdfsubject={HAAP SDK v0.1.0-alpha.11},
]{hyperref}

% ── Headers & footers ────────────────────────────────────────────────────────
\usepackage{fancyhdr}
\pagestyle{fancy}
\fancyhf{}
\fancyhead[L]{\small\color{codegray}HAAP SDK — Distribution \& CrewAI Guide}
\fancyhead[R]{\small\color{codegray}v0.1.0-alpha.11}
\fancyfoot[C]{\small\thepage}
\renewcommand{\headrulewidth}{0.4pt}
\renewcommand{\headrule}{\color{frameborder}\hrule width\headwidth height\headrulewidth}

% ── Section title formatting ─────────────────────────────────────────────────
\usepackage{titlesec}
\titleformat{\section}
  {\large\bfseries\color{hawcxblue}}{\thesection}{0.8em}{}[\vspace{-4pt}\color{frameborder}\hrule\vspace{2pt}]
\titleformat{\subsection}
  {\normalsize\bfseries}{\thesubsection}{0.6em}{}
\titleformat{\subsubsection}
  {\normalsize\itshape}{\thesubsubsection}{0.6em}{}

% ── Misc utilities ───────────────────────────────────────────────────────────
\usepackage{enumitem}
\setlist[itemize]{leftmargin=1.5em, itemsep=2pt, topsep=4pt}
\setlist[enumerate]{leftmargin=1.5em, itemsep=2pt, topsep=4pt}

\newcommand{\pkg}[1]{\texttt{#1}}
\newcommand{\code}[1]{\texttt{#1}}

% ─────────────────────────────────────────────────────────────────────────────

\begin{document}

% ── Title ────────────────────────────────────────────────────────────────────
\begin{center}
  {\LARGE\bfseries\color{hawcxblue} HAAP SDK}\\[6pt]
  {\large Binary Distribution \& CrewAI Integration Guide}\\[10pt]
  {\small SDK version \texttt{0.1.0-alpha.11} \quad|\quad
   Protocol: HAAP Canonical Specification v7.2.5 \quad|\quad
   2026-05-26}\\[4pt]
  {\small\color{codered}\itshape Pre-product alpha. Customer evaluation only.}
\end{center}

\vspace{4pt}
\tableofcontents
\vspace{8pt}

% ── 1. What Is Distributed ───────────────────────────────────────────────────
\section{What Is Distributed}

The HAAP SDK ships \textbf{one real binary} --- \pkg{hawcx-manager} ---
packaged into platform-specific wheels (PyPI) and scoped npm packages (npm).
The binary implements the complete customer-side pipeline described in HAAP
CS~v7.2.5~§39 (Profile~E, five-process pipeline).

\begin{table}[h!]
\centering
\caption{Components dispatched by \texttt{hawcx-manager}}
\renewcommand{\arraystretch}{1.25}
\begin{tabularx}{\linewidth}{@{} l l X @{}}
\toprule
\textbf{Role} & \textbf{Subcommand / Legacy name} & \textbf{Spec reference} \\
\midrule
Supervisor          & \code{supervisor} / \code{haap-supervisor}           & §39.1 — lifecycle orchestrator \\
Authenticator       & \code{authenticator} / \code{haap-authenticator}     & §4.2.1 — holds \code{IK\_i}, runs X3DH \\
TQS Precompute      & \code{tqs-precompute} / \code{haap-tqs-precompute}   & §40 — session-scoped \code{K\_session\_root} \\
TQS JIT             & \code{tqs-jit} / \code{haap-tqs-jit}                 & §40 — request-scoped, 1:1 with Assembler \\
Assembler           & \code{assembler} / \code{haap-assembler}             & §39, §47 — single-flight crypto-proxy \\
EIB                 & \code{eib} / \code{haap-eib}                         & §45 — External Identity Broker, OAuth bearer \\
SDK CLI             & \code{sdk} / \code{haap-sdk}                         & debug / operator CLI \\
\bottomrule
\end{tabularx}
\end{table}

Seven legacy names (\pkg{haap-supervisor}, \pkg{haap-authenticator}, etc.)\ are
preserved as symlinks (Unix) or \texttt{.exe} copies (Windows) so existing
deployment scripts continue to work unchanged.

\begin{mdframed}[style=notebox]
\textbf{Not in this package:} \pkg{haap-rsv} (MCP-server-side verifier) ships
from a separate image \texttt{ghcr.io/hawcx/haap-rsv} built from
\texttt{hx\_agent\_authorizer}. It is Unix-only by design (UDS + peer-credential
checks per §39.12.1).
\end{mdframed}

% ── 2. PyPI Distribution ─────────────────────────────────────────────────────
\section{PyPI Distribution (\texttt{hawcx-haap})}

\subsection{Install}

\begin{lstlisting}[language=bash, numbers=none, frame=none, backgroundcolor=\color{codebg}]
pip install hawcx-haap
\end{lstlisting}

On \code{pip install}, pip selects the wheel matching the current platform.
\pkg{hawcx-manager} lands in \code{venv/bin/hawcx-manager} (Unix) or
\code{venv\textbackslash Scripts\textbackslash hawcx-manager.exe} (Windows).

\subsection{Build backend: maturin \texttt{bindings = "bin"}}

The build backend is \href{https://www.maturin.rs/}{maturin} configured with
\code{bindings = "bin"}. This packages a Rust binary --- not a Python
extension --- into a platform-specific wheel. The same mechanism is used by
\pkg{ruff}, \pkg{uv}, and \pkg{black}.

\begin{lstlisting}[language=TOML, caption={python/pyproject.toml (relevant section)}]
[build-system]
requires = ["maturin>=1.5,<2.0"]
build-backend = "maturin"

[tool.maturin]
bindings = "bin"
# Rust source lives in sibling hx_agent_client_auth_service repo.
# CI passes --manifest-path at build time via release-python.yml.
python-source = "src"
\end{lstlisting}

\subsection{Supported platforms}

\begin{table}[h!]
\centering
\caption{PyPI wheel platform tags}
\renewcommand{\arraystretch}{1.2}
\begin{tabular}{@{} l l l @{}}
\toprule
\textbf{Wheel tag suffix} & \textbf{OS} & \textbf{Architecture} \\
\midrule
\code{linux\_x86\_64}           & Linux   & x86-64  \\
\code{manylinux\_*\_aarch64}    & Linux   & ARM64   \\
\code{macosx\_*\_arm64}         & macOS   & ARM64 (Apple Silicon) \\
\code{win\_amd64}               & Windows & x86-64  \\
\code{win\_arm64}               & Windows & ARM64   \\
\bottomrule
\end{tabular}
\end{table}

\subsection{Locating the binary at runtime}

\begin{lstlisting}[language=Python, caption={Resolving the bundled binary path}]
from hawcx_haap import get_binary_path

binary = get_binary_path()
# Unix:    "/path/to/venv/bin/hawcx-manager"
# Windows: "C:\...\venv\Scripts\hawcx-manager.exe"
\end{lstlisting}

\code{get\_binary\_path()} raises \code{RuntimeError} if the binary is absent
(e.g., installed from source without maturin). The error message explains the
local development workflow.

\subsection{PyPI packages}

\begin{table}[h!]
\centering
\renewcommand{\arraystretch}{1.2}
\begin{tabular}{@{} l X @{}}
\toprule
\textbf{Package} & \textbf{Role} \\
\midrule
\pkg{hawcx-haap}   & Core IPC client + bundled \pkg{hawcx-manager} binary \\
\pkg{hawcx-crewai} & CrewAI \code{BaseTool} adapter (depends on \pkg{hawcx-haap}) \\
\bottomrule
\end{tabular}
\end{table}

% ── 3. npm Distribution ──────────────────────────────────────────────────────
\section{npm Distribution (\texttt{@hawcx/hawcx-haap})}

\subsection{Install}

\begin{lstlisting}[language=bash, numbers=none, frame=none, backgroundcolor=\color{codebg}]
npm install @hawcx/hawcx-haap
\end{lstlisting}

npm installs only the platform package matching the current OS and CPU.

\subsection{optionalDependencies pattern}

The main package declares \code{optionalDependencies} pointing at five scoped
packages, each tagged with \code{os} and \code{cpu} fields. npm skips packages
for non-matching platforms. The pattern is identical to \pkg{@biomejs/biome}
and \pkg{esbuild}.

\begin{lstlisting}[language=JSON, caption={node/package.json (optionalDependencies)}]
"optionalDependencies": {
  "@hawcx/hawcx-haap-linux-x64":    "0.1.0-alpha.11",
  "@hawcx/hawcx-haap-linux-arm64":  "0.1.0-alpha.11",
  "@hawcx/hawcx-haap-darwin-arm64": "0.1.0-alpha.11",
  "@hawcx/hawcx-haap-win32-x64":    "0.1.0-alpha.11",
  "@hawcx/hawcx-haap-win32-arm64":  "0.1.0-alpha.11"
}
\end{lstlisting}

\subsection{Supported platforms}

\begin{table}[h!]
\centering
\caption{npm platform packages}
\renewcommand{\arraystretch}{1.2}
\begin{tabular}{@{} l l l @{}}
\toprule
\textbf{npm package} & \textbf{OS} & \textbf{CPU} \\
\midrule
\code{@hawcx/hawcx-haap-linux-x64}    & linux  & x64   \\
\code{@hawcx/hawcx-haap-linux-arm64}  & linux  & arm64 \\
\code{@hawcx/hawcx-haap-darwin-arm64} & darwin & arm64 \\
\code{@hawcx/hawcx-haap-win32-x64}    & win32  & x64   \\
\code{@hawcx/hawcx-haap-win32-arm64}  & win32  & arm64 \\
\bottomrule
\end{tabular}
\end{table}

\subsection{Locating the binary at runtime}

\begin{lstlisting}[language={}  , caption={TypeScript --- resolving the bundled binary}]
import { getBinaryPath } from "@hawcx/hawcx-haap";

const binary = getBinaryPath();
// -> "/path/to/.../node_modules/@hawcx/hawcx-haap-linux-x64/hawcx-manager"
\end{lstlisting}

% ── 4. CI Build Matrix ───────────────────────────────────────────────────────
\section{CI Build Matrix}

Both \code{release-python.yml} and \code{release-node.yml} run the same
five-target Rust build matrix. The binary is checked out from
\code{hx\_agent\_client\_auth\_service} and built against the private Kellnr
registry at \code{cargo.hawcx.com}.

\begin{table}[h!]
\centering
\caption{Rust build targets and CI runners}
\renewcommand{\arraystretch}{1.25}
\begin{tabular}{@{} l l l l @{}}
\toprule
\textbf{Rust target} & \textbf{CI runner} & \textbf{PyPI tag} & \textbf{npm pkg} \\
\midrule
\code{x86\_64-unknown-linux-gnu}   & \code{ubuntu-22.04}     & \code{linux\_x86\_64}        & \code{linux-x64}    \\
\code{aarch64-unknown-linux-gnu}   & \code{ubuntu-22.04-arm} & \code{*\_aarch64}            & \code{linux-arm64}  \\
\code{aarch64-apple-darwin}        & \code{macos-14}         & \code{macosx\_*\_arm64}      & \code{darwin-arm64} \\
\code{x86\_64-pc-windows-msvc}     & \code{windows-latest}   & \code{win\_amd64}            & \code{win32-x64}    \\
\code{aarch64-pc-windows-msvc}     & \code{windows-latest}   & \code{win\_arm64}            & \code{win32-arm64}  \\
\bottomrule
\end{tabular}
\end{table}

Releases are triggered by pushing version tags:

\begin{lstlisting}[language=bash, numbers=none, frame=none, backgroundcolor=\color{codebg}]
git tag python-v0.1.0-alpha.11 && git push origin python-v0.1.0-alpha.11
git tag node-v0.1.0-alpha.11   && git push origin node-v0.1.0-alpha.11
\end{lstlisting}

% ── 5. CrewAI Integration ────────────────────────────────────────────────────
\section{Getting Started with CrewAI}

\subsection{Prerequisites}

\begin{enumerate}
  \item The \pkg{hawcx-manager} supervisor pipeline must be running on the
        agent host. The Assembler's socket follows the convention:\\
        \code{\$XDG\_RUNTIME\_DIR/hawcx/\{agent\_id\}/agent-assembler-0.sock}
        (Linux) or\\
        \code{\textbackslash\textbackslash.\textbackslash pipe\textbackslash haap-\{agent\_id\}-agent-assembler-0}
        (Windows).
  \item The agent identity must be pre-provisioned via the Hawcx Admin Console
        (CAA \textrightarrow\ Authenticator flow per CS~§4.6.3).
\end{enumerate}

\subsection{Install}

\begin{lstlisting}[language=bash, numbers=none, frame=none, backgroundcolor=\color{codebg}]
pip install hawcx-haap hawcx-crewai crewai
\end{lstlisting}

\subsection{Package roles}

\begin{table}[h!]
\centering
\renewcommand{\arraystretch}{1.2}
\begin{tabular}{@{} l X @{}}
\toprule
\textbf{Package} & \textbf{Role} \\
\midrule
\pkg{hawcx-haap}   & IPC client + bundled binary. Core SDK. \\
\pkg{hawcx-crewai} & \code{HawcxTool} --- a \code{crewai.tools.BaseTool} subclass. \\
\pkg{crewai}       & Multi-agent orchestration framework. \\
\bottomrule
\end{tabular}
\end{table}

% ── 6. Annotated CrewAI Example ──────────────────────────────────────────────
\section{Annotated CrewAI Example}

The following end-to-end example shows the idiomatic pattern for a
HAAP-authenticated CrewAI crew. Annotations explain every security-relevant
decision.

\subsubsection*{Step 1 — Define typed argument schemas}

CrewAI surfaces the \code{args\_schema} to the LLM for argument validation and
tool-selection prompting. The \code{user\_principal\_id} field is the per-call
identity axis: the LLM supplies it, and the SDK enforces that it is in the
operator-controlled allowlist before writing any IPC bytes (HAAP CS~v7.2.5
H-3 hardening).

\begin{lstlisting}[language=Python, caption={Pydantic schemas for tool arguments}]
from pydantic import BaseModel, Field

class SearchInput(BaseModel):
    query: str = Field(
        description="Search query against the NIM knowledge base."
    )
    user_principal_id: str = Field(
        description=(
            "End-user on whose behalf the search runs. "
            "Must be a principal registered for this agent."
        )
    )

class DocsInput(BaseModel):
    document_id: str = Field(
        description="Opaque document ID to retrieve."
    )
    user_principal_id: str = Field(
        description="End-user on whose behalf the document is fetched."
    )
\end{lstlisting}

\subsubsection*{Step 2 --- Connect to the Assembler}

\code{HawcxAgent.connect\_by\_agent\_id()} resolves the conventional UDS path
and performs the §7 capability handshake synchronously. \code{principal\_allowlist}
is a required, \emph{operator-controlled} closed set. Never derive it from LLM
output or request bodies.

\begin{lstlisting}[language=Python, caption={Opening the Assembler IPC connection}]
from hawcx_haap import HawcxAgent

ALLOWED_USERS = ["alice@example.com", "bob@example.com"]

with HawcxAgent.connect_by_agent_id(
    "research-u1",
    principal_allowlist=ALLOWED_USERS,  # required; enforced before any IPC write
) as agent:
    ...
\end{lstlisting}

\subsubsection*{Step 3 --- Construct \texttt{HawcxTool} instances}

One \code{HawcxTool} per logical tool. All instances share the same
\code{agent} (one Assembler connection per process). Construction is cheap ---
no IPC at this point.

Key fields:

\begin{itemize}
  \item \textbf{\code{provider}} --- §47.2 provider class. Routes the §47.8
        \code{GetExternalCredential} IPC to the correct sidecar
        (\code{haap-nim-provider}, \code{haap-anthropic-provider}, etc.).
  \item \textbf{\code{tool\_id}} --- §47.4 tool identity binding. The sidecar
        refuses to disclose a credential unless the requesting \code{tool\_id}
        matches the bound value. Two tools with different \code{tool\_id}s
        cannot share credentials even within the same pipeline.
  \item \textbf{\code{endpoint}} --- Destination URL. The Assembler constructs
        the outbound HTTPS request; the Python process never sees the bearer
        token.
  \item \textbf{\code{action}} --- TBAC action list for Cedar policy evaluation.
\end{itemize}

\begin{lstlisting}[language=Python, caption={Constructing HawcxTool instances}]
from hawcx_crewai import HawcxTool

nim_search_tool = HawcxTool(
    name="nim_search",
    description=(
        "Search the organisation's private knowledge base via NVIDIA NIM. "
        "Always supply the user_principal_id you were given."
    ),
    hawcx_agent=agent,
    provider="nim",                          # routes to haap-nim-provider sidecar
    tool_id="nim-search-v1",                 # credential binding key (sec. 47.4)
    endpoint="https://api.nim.nvidia.com/v1/search",
    method="POST",
    action=["read"],                         # Cedar policy action
    args_schema=SearchInput,
)

docs_tool = HawcxTool(
    name="docs_reader",
    description=(
        "Retrieve a document from the internal store by ID. "
        "Always supply the user_principal_id you were given."
    ),
    hawcx_agent=agent,
    provider="generic-bearer",               # generic bearer sidecar
    tool_id="docs-reader-v1",               # separate credential from nim_search
    endpoint="https://docs.internal.example.com/v1/documents",
    method="GET",
    action=["read"],
    args_schema=DocsInput,
)
\end{lstlisting}

\subsubsection*{Step 4 --- Define CrewAI Agents}

Each CrewAI \code{Agent} receives the relevant \code{HawcxTool}(s) in its
\code{tools} list. The LLM backing the agent sees only the tool's name,
description, and \code{args\_schema} --- never provider keys, HAAP session
tokens, or IPC internals.

\begin{lstlisting}[language=Python, caption={Defining CrewAI agents}]
from crewai import Agent

researcher = Agent(
    role="Research Analyst",
    goal=(
        "Find and retrieve the most relevant technical documents "
        "answering the user's question."
    ),
    backstory=(
        "A meticulous analyst with access to a private knowledge base. "
        "Always cites sources by document ID."
    ),
    tools=[nim_search_tool, docs_tool],
    verbose=True,
)

summarizer = Agent(
    role="Technical Writer",
    goal="Synthesize research findings into a concise technical summary.",
    backstory=(
        "Turns raw research into polished summaries. "
        "Does not call external tools."
    ),
    tools=[],          # no HAAP tools needed for summarisation
    verbose=True,
)
\end{lstlisting}

\subsubsection*{Step 5 --- Define Tasks}

The \code{user\_principal\_id} is injected into the task description so the
LLM learns to pass it through every tool call. The principal is sourced from
the authenticated session --- never from LLM output.

\begin{lstlisting}[language=Python, caption={Task definitions with per-call principal}]
from crewai import Task

user_id  = "alice@example.com"    # from authenticated session
question = "What are the performance characteristics of the HAAP TQS pipeline?"

research_task = Task(
    description=(
        f"Answer: {question}\n\n"
        f"IMPORTANT: For every tool call, pass "
        f"user_principal_id='{user_id}'. Do not omit this field."
    ),
    expected_output=(
        "List of relevant document IDs with key findings and direct quotes."
    ),
    agent=researcher,
)

summarize_task = Task(
    description=(
        "Using the research findings above, write a 3-paragraph technical "
        "summary for a developer audience. Include document IDs as citations."
    ),
    expected_output="Concise 3-paragraph technical summary with citations.",
    agent=summarizer,
    context=[research_task],      # receives researcher's output as context
)
\end{lstlisting}

\subsubsection*{Step 6 --- Assemble and run the Crew}

\begin{lstlisting}[language=Python, caption={Assembling and running the Crew}]
from crewai import Crew, Process

crew = Crew(
    agents=[researcher, summarizer],
    tasks=[research_task, summarize_task],
    process=Process.sequential,   # tasks run in order; use hierarchical
    verbose=True,                  # for manager-agent architecture
)

result = crew.kickoff()
print(result)
# The `with` block closes the Assembler IPC socket on exit.
# Keep the agent open for the lifetime of the process to avoid
# the ~5 ms per-reconnect handshake latency.
\end{lstlisting}

% ── 7. Multi-Tenant Pattern ──────────────────────────────────────────────────
\section{Multi-Tenant Pattern}

When one agent process serves many end-users, two patterns are supported:

\subsubsection*{A. Per-call principal via \texttt{args\_schema} (recommended)}

Declare \code{user\_principal\_id} in the tool's \code{args\_schema}. The LLM
supplies it on each call; the SDK enforces it against \code{principal\_allowlist}
before any IPC.

\begin{lstlisting}[language=Python, caption={Per-call principal in args\_schema}]
# SearchInput / DocsInput in Step 1 already include user_principal_id.
# No extra code is needed; the SDK extracts it from kwargs in HawcxTool._run.
\end{lstlisting}

\subsubsection*{B. Per-user tool instance via \texttt{for\_user}}

For flows that materialise one tool per user request:

\begin{lstlisting}[language=Python, caption={for\_user sugar}]
# Cheap -- shares the underlying HawcxAgent (one IPC connection).
alice_search = nim_search_tool.for_user("alice@example.com")
bob_search   = nim_search_tool.for_user("bob@example.com")

# alice_search always calls invoke_for("alice@example.com", ...)
# bob_search   always calls invoke_for("bob@example.com",   ...)
\end{lstlisting}

% ── 8. Security Model ────────────────────────────────────────────────────────
\section{Security Model}

\begin{table}[h!]
\centering
\caption{HAAP SDK security properties}
\renewcommand{\arraystretch}{1.35}
\begin{tabularx}{\linewidth}{@{} p{4.5cm} X @{}}
\toprule
\textbf{Property} & \textbf{How it is achieved} \\
\midrule
Credentials never reach the LLM process &
  \code{HawcxTool.\_run} calls \code{HawcxAgent.invoke} over local UDS only.
  The Assembler fetches provider credentials internally (§47.8
  \code{GetExternalCredential}) and attaches them to the outbound HTTPS
  request. The credential value never returns to Python. \\
\addlinespace
HAAP session keys never reach the LLM process &
  All crypto (\code{K\_session\_root}, \code{K\_req}, \code{K\_resp}) lives
  inside the Assembler process. The SDK exchanges only plaintext request
  bodies and decrypted response bodies. \\
\addlinespace
LLM-supplied principals are sandboxed &
  \code{principal\_allowlist} is a closed, operator-controlled set validated
  synchronously at \code{HawcxAgent.connect()} time. A compromised or
  hallucinating LLM cannot escalate to an out-of-list principal. \\
\addlinespace
Per-tool credential binding &
  \code{tool\_id} (§47.4) prevents one tool from borrowing another tool's
  credentials even within the same pipeline. \\
\addlinespace
IPC socket isolation &
  Sockets are placed under \code{\$XDG\_RUNTIME\_DIR/hawcx/} (0o700, per-UID,
  systemd-managed). The SDK refuses to fall back to
  \code{/tmp/hawcx/} without an explicit
  \code{HAAP\_SDK\_ALLOW\_TMP\_IPC=1} opt-in. \\
\bottomrule
\end{tabularx}
\end{table}

\vspace{8pt}
\noindent\rule{\linewidth}{0.4pt}\\[4pt]
{\small\color{codegray}
Generated 2026-05-26. Source: \texttt{hawcx\_agentic\_sdk} \texttt{0.1.0-alpha.11},
HAAP Canonical Specification v7.2.5. Copyright \textcopyright\ 2026 Hawcx Inc.
All rights reserved.}

\end{document}
